\documentclass[11pt,letterpaper]{article}
\usepackage[letterpaper,left=0.85in,right=0.85in,top=0.78in,bottom=0.80in,headheight=14pt]{geometry}
\usepackage[T1]{fontenc}
\usepackage{lmodern}
\usepackage{amsmath,amssymb,mathtools}
\usepackage{microtype}
\usepackage{array,tabularx,booktabs}
\usepackage{xcolor}
\usepackage{tikz}
\usepackage[most]{tcolorbox}
\usepackage{titlesec,fancyhdr}
\usepackage{hyperref}
\definecolor{ink}{HTML}{172126}
\definecolor{navy}{HTML}{17324D}
\definecolor{teal}{HTML}{1A7F86}
\definecolor{tealdeep}{HTML}{155E63}
\definecolor{mist}{HTML}{F3F8F8}
\definecolor{linegray}{HTML}{CFDADB}
\definecolor{linkgray}{HTML}{315C66}
\hypersetup{colorlinks=true,urlcolor=linkgray,linkcolor=ink,citecolor=ink,
pdftitle={The Not-So-Hard Problem of Consciousness},pdfauthor={Michael Zot},
pdfsubject={Empirically constrained theoretical and methodological paper}}
\setlength{\parindent}{0pt}
\setlength{\parskip}{0.48em}
\setlength{\emergencystretch}{1.8em}
\linespread{1.06}
\clubpenalty=10000
\widowpenalty=10000
\titleformat{\section}{\sffamily\large\bfseries\color{navy}}{\thesection.}{0.55em}{}
\titleformat{\subsection}{\sffamily\normalsize\bfseries\color{tealdeep}}{}{0pt}{}
\titlespacing*{\section}{0pt}{1.15em}{0.4em}
\titlespacing*{\subsection}{0pt}{0.7em}{0.2em}
\pagestyle{fancy}\fancyhf{}
\fancyhead[L]{\scriptsize\sffamily Michael Zot / The Not-So-Hard Problem of Consciousness}
\fancyfoot[L]{\scriptsize\sffamily SCTC | Theoretical \& Methodological Paper | 1 October 2026}
\fancyfoot[R]{\scriptsize\sffamily\thepage}
\renewcommand{\headrulewidth}{0.25pt}
\fancypagestyle{firstpage}{\fancyhf{}
\fancyfoot[L]{\scriptsize\sffamily SCTC | Theoretical \& Methodological Paper | 1 October 2026}
\fancyfoot[R]{\scriptsize\sffamily\thepage}
\renewcommand{\headrulewidth}{0pt}}
\newcommand{\Csys}{C_{\mathrm{sys},B}(t)}
\newcommand{\Lsys}{L_B(t)}
\newcommand{\repoURL}{https://github.com/mikecreation/ZotBot/tree/main/research/not-so-hard-problem-of-consciousness}
\newtcolorbox{abstractcard}{
  enhanced, breakable, colback=white, colframe=linegray,
  boxrule=0.8pt, arc=2.5mm, left=4mm, right=4mm, top=3mm, bottom=3mm,
  borderline west={2.2pt}{0pt}{teal},
  title={\sffamily\bfseries\color{white} Abstract},
  colbacktitle=navy, coltitle=white,
  fonttitle=\normalsize, attach boxed title to top left={xshift=2mm,yshift=-2mm},
  boxed title style={arc=1.5mm,boxrule=0pt,left=2.7mm,right=2.7mm,top=1.1mm,bottom=1.1mm}
}
\newtcolorbox{keybox}[1]{
  enhanced, breakable, colback=mist, colframe=linegray,
  boxrule=0.6pt, arc=2mm, left=3.5mm, right=3.5mm, top=2.7mm, bottom=2.7mm,
  borderline west={2pt}{0pt}{teal},
  title={\sffamily\bfseries #1}, colbacktitle=mist, coltitle=navy,
  boxed title style={boxrule=0pt,colback=mist,left=0pt,right=0pt,top=0pt,bottom=0pt},
  attach title to upper={\par\smallskip}
}
\begin{document}
\thispagestyle{firstpage}

\begin{tcolorbox}[
  enhanced, boxrule=0pt, arc=3mm,
  interior style={left color=navy,right color=teal},
  left=6mm,right=6mm,top=5mm,bottom=5mm
]
{\scriptsize\sffamily\bfseries\color{white!88} THEORETICAL \& METHODOLOGICAL PAPER \hfill 1 OCTOBER 2026\par}
\vspace{1.0em}
{\fontsize{30}{33}\selectfont\bfseries\color{white} The Not-So-Hard Problem\par of Consciousness\par}
\vspace{0.8em}
{\large\sffamily\color{white!94} Defining the conscious system before explaining its changing states\par}
\end{tcolorbox}

\vspace{0.65em}
{\normalsize\textbf{Michael Zot} \quad \textcolor{linegray}{|} \quad Independent Researcher, Illinois, USA\par}
{\small \href{mailto:mike@ZotBot.Ai}{mike@ZotBot.Ai}\par}
{\small\sffamily\textbf{Research repository:} \href{\repoURL}{github.com/mikecreation/ZotBot/tree/main/research/not-so-hard-problem-of-consciousness}\par}
\vspace{0.55em}

\begin{abstractcard}
\small
\textbf{A biological individual can continue while its responsiveness, memory, wakefulness, and access to the environment change.} The persistence of that individual, the occurrence of an experience, and our ability to detect the experience therefore require separate scientific descriptions. This paper develops the \textbf{Systemic Continuity Theory of Consciousness (SCTC)}: a definition of minimal system-level consciousness as the viable existence of an integrated biological individual, tracked through its organismal continuity. Phenomenal experience is represented separately as $E$.

SCTC is evaluated against existing empirical dissociations. Neural command following without an observable command response, dreaming during sleep, and severe memory impairment with other abilities preserved constrain the properties that can define the continuing biological system. SCTC makes its system-level assignment explicitly; the empirical literature independently supports the separations on which the framework operates.

The argument distinguishes a persistence condition from a changing faculty and an observed signal from the state inferred through it. A faculty that can fail while the individual continues cannot be a necessary condition of that individual's persistence. This excludes several familiar measures as definitions of the system-level target. Biological continuity supplies a minimal account of that target without making experience, memory, or report interchangeable with it. Rival definitions are compared by empirical fit, minimality, boundary consistency, and whether they address the same object.

The result is a definition grounded in an explicit target and constrained by published evidence. No new experiment is required to establish dissociations already observed. The further claim that this definition uniquely captures every legitimate meaning of consciousness does not follow from those dissociations. The hard question of why experience occurs remains, but it concerns $E$; it does not settle whether the biological individual persists.
\end{abstractcard}

\begin{keybox}{What this paper separates}
\small
\textbf{The continuing biological individual} is tracked separately from \textbf{experience, memory, wakefulness, responsiveness, cognitive access, report, and the evidence used to infer those states}. Evidence for one variable cannot automatically stand in for evidence for another.
\end{keybox}

{\small\textbf{Article type:} Empirically constrained theoretical and methodological paper. Published studies provide the empirical basis; no new experimental dataset is reported.\par}
{\small\textbf{Keywords:} consciousness; biological continuity; phenomenal experience; definition; neurological dissociation; scientific inference.\par}
\vspace{0.25em}
\noindent\textcolor{linegray}{\rule{\linewidth}{0.45pt}}
\vspace{0.25em}

\section{Define the object before explaining its states}
Imagine waking after a procedure and remembering nothing. Did the biological individual continue? Did any experience occur? Was any experience retained so it could later be recalled? These are three questions. An answer to the last does not automatically answer the other two.

The distinction becomes clear in a dream. The sleeper continues as a biological individual; a particular experience may occur; its later recall may succeed or fail. The dream is a state of the individual. Remembering it is another event. The evidence available afterward is a third. A scientific account must say which of these it defines and which it explains.

SCTC places the continuing individual at the system level and its mental states at the faculty level. Its central claim is that \textbf{a variable used to define the persistence of the system must not disappear merely because one of the system's faculties or observation channels fails}. This is a requirement on a persistence-based target. It is not a requirement that every experiential state remain constant while an organism lives.

Target separation has a substantial history. Chalmers distinguishes experience from functional problems; Block distinguishes phenomenal from access consciousness; Bayne and colleagues distinguish multiple dimensions of global states \cite{Chalmers1995,Block1995,Bayne2016}. SCTC uses these distinctions while giving biological persistence an explicit system-level role. Its contribution is the argument for that role and the constraints it imposes on measurement and definition.

\section{The definition and its necessary constraints}
\subsection*{System, state, and continuity}
Let $B$ denote the biological individual under study. Let $L_B(t)=1$ mean that $B$ exists as a viable, integrated biological individual at time $t$. Let $\Gamma_B(t_0,t_1)=1$ mean that the individual at $t_1$ continues the organismal causal trajectory of the individual at $t_0$. A classification is left unresolved when the evidence does not settle the relevant biological boundary.

$L$ and $\Gamma$ answer different questions. Two individuals can both be alive without being the same individual. Constancy of $L$ is therefore not sufficient to establish $\Gamma$. Continuity requires causal history as well as a defensible account of biological individuality.

SCTC defines its minimal system-level target by
\begin{equation}
\Csys\coloneqq\Lsys.
\label{eq:definition}
\end{equation}
In plain language: \textbf{system-level consciousness is the viable existence of the biological individual; continuity tracks that individual across changes in its states.} The qualifier \emph{system-level} remains explicit throughout the argument.

Phenomenal experience, $E$, means that there is something it feels like to be in a state: seeing red, feeling pain, or dreaming. Cognitive access, $Q$, means that information is available for a specified use. Behavioral response, $R$, and later report, $H$, are outputs. Memory, $M$, concerns retention and retrieval. These variables may interact and share mechanisms. Giving them separate symbols makes their relations questions to investigate rather than identities to assume.

\subsection*{The faculty-exclusion argument}
Suppose a candidate definition says that faculty $F$ is necessary for biological persistence. A reliably established case in which the individual persists while that faculty is absent defeats the claim of necessity:
\begin{equation}
\exists(B,t):\quad L_B(t)=1\ \land\ F_B(t)=0
\quad\Longrightarrow\quad
F\text{ is not necessary for }L.
\label{eq:exclusion}
\end{equation}
This is ordinary counterexample reasoning. The experiment need not be new: an already published counterexample still excludes a universal necessity claim. What matters is whether the observation establishes the relevant absence and persistence at the specified scale.

Under the definition in Eq.~\ref{eq:definition}, the exclusion applies to $C_{\mathrm{sys}}$ as well. The logical order matters. The evidence first separates a faculty from biological persistence. SCTC then identifies persistence as its system-level target. Calling the continuing organism a \emph{conscious individual} cannot independently prove that assignment.

The same evidence does not exclude the faculty from every other target. An experience-centered account can consistently hold that an organism persists during an interval without experience. It is defining an experiential condition, whereas SCTC is defining organismal existence. Comparing them requires keeping that difference visible.

\section{Existing evidence constrains candidate definitions}
\subsection*{Responsiveness and cognitive command following}
Bodien and colleagues assessed 353 adults with disorders of consciousness. Of the 241 participants without an observable response to commands, 60 (25\%) showed task-based command following on fMRI or EEG \cite{Bodien2024}. The percentage refers to this subgroup in a convenience sample. It is not an estimate for every unresponsive patient.

The result excludes overt command response as a necessary indicator of the detected command-related cognition. It also shows why a failed motor channel cannot, by itself, define the disappearance of the biological individual or all cognitive capacity. The neural task response supplies additional evidence; it does not reveal every feature of experience or establish its absence in participants with negative tests.

\subsection*{Wakefulness, environmental connection, and experience}
Sanders and colleagues distinguish experience, connectedness to the environment, and responsiveness \cite{Sanders2012}. Siclari and colleagues relate sleep EEG patterns to reports obtained after awakenings \cite{Siclari2017}. Dreaming illustrates the separation between ordinary wakefulness, engagement with the surrounding room, and an experienced state.

These findings constrain definitions that treat wakefulness or overt response as interchangeable with experience. They also support tracking the same individual across waking and sleeping states. Recall still mediates the evidence for particular dreams; the presence of sleep alone establishes neither the presence nor absence of experience.

\subsection*{Memory and the earlier occurrence of a state}
Medial temporal-lobe amnesia can profoundly impair the acquisition of new memories while other abilities, including perception, language, and reasoning, remain relatively preserved \cite{Milner2005}. This is a dissociation involving particular memory functions. It does not establish the disappearance of every memory system.

The consequence is precise: the ability to recall an earlier state cannot simply be identified with the occurrence of that state or with persistence of the individual. Later memory is one route to evidence about earlier experience. Failure of that route requires an account of the memory process before it can support a stronger verdict.

\subsection*{A measured signature must keep its measured target}
Casey and colleagues evaluate EEG signatures using explicit definitions of responsiveness, connectedness, and experience \cite{Casey2024}. Casali and colleagues develop the perturbational complexity index using TMS-EEG \cite{Casali2013}. Such work demonstrates that consciousness research already uses distinctions and evidence beyond overt behavior.

The SCTC constraint is to retain those distinctions when interpreting a result. A detector calibrated to a response contrast cannot automatically inherit sensitivity to all experience. Conversely, the absence of overt behavior does not make instrumental evidence irrelevant. Each channel must be evaluated against the target it actually supports.

\subsection*{What the accumulated record establishes}
\begin{center}
\small
\renewcommand{\arraystretch}{1.2}
\begin{tabularx}{\linewidth}{@{}>{\raggedright\arraybackslash}p{0.24\linewidth} X@{}}
\toprule
\textbf{Candidate property} & \textbf{Constraint from the cited evidence}\\
\midrule
Overt command response & Can be absent despite detectable command-related cognition.\\
Ordinary wakefulness & Does not exhaust the conditions in which dream experience is reported.\\
Episodic memory formation & Can be profoundly impaired while the individual and other abilities persist.\\
Later report or recall & Depends on channels distinct from the earlier state being inferred.\\
Phenomenal experience & Must be assessed separately; organismal persistence alone does not establish its occurrence.\\
\bottomrule
\end{tabularx}
\end{center}
These constraints already exist. Their use here is empirical evaluation of a definition, not an announcement of a new experiment. Pairwise dissociations do not establish that every faculty can be absent simultaneously. They establish the particular separations documented in the relevant studies.

\section{The invariant and the limits of inference}
\subsection*{A stable referent across changing configurations}
In the cited clinical and sleep contexts, researchers follow a biological individual whose state or detectable performance changes. SCTC assigns the system-level referent to that individual, rather than letting a change in one faculty redefine the biological unit under study.

This yields a useful distinction between \emph{what persists} and \emph{what varies}. It does not assert that the entire organism remains physically unchanged. Metabolism, repair, injury, and support alter the implementation. The continuity claim concerns the organismal trajectory through those changes, subject to independent criteria for individuality.

Biological continuity is the minimal candidate for this declared persistence target because it does not require an additional experiential or behavioral condition to identify the persisting individual. Other formulations can describe the same trajectory. The minimality claim is about the target specified here; it does not establish a unique definition of every phenomenon historically called consciousness.

\subsection*{Why a missing response cannot locate the failure}
Consider an ideal model of an immediate positive report of experience. Let $E$ indicate experience, $A$ an available access route for the task, and $O$ an operational output channel. Assume that the report $R_{\mathrm{rep}}$ occurs exactly when all three binary factors are present:
\begin{equation}
R_{\mathrm{rep}}=EAO.
\label{eq:report}
\end{equation}
If $R_{\mathrm{rep}}=0$, experience may be absent, access may fail, or output may be blocked. The same observation fits different hidden conditions. This is \emph{non-identifiability}: the observation alone does not specify which cause produced it. Actual reports also involve comprehension, attention, effort, and error; Eq.~\ref{eq:report} is a transparent counterexample to an automatic inference, not a complete theory of behavior.

For later recall, retention and retrieval add further dependencies. The individual can persist throughout while an experience occurs, fails to occur, or remains unresolved by the available evidence. SCTC keeps those alternatives separate.

\subsection*{Negative evidence still has a role}
A negative result can lower the probability of experience without proving absence. Let $q$ be the initial probability of experience, $s$ the chance of a positive test with experience, and $c$ the chance of a negative test without it. Bayes' rule gives
\begin{equation}
P(E=1\mid D=0)=\frac{q(1-s)}{q(1-s)+(1-q)c}.
\label{eq:bayes}
\end{equation}
For $0<q<1$ and a defined denominator, the updated probability decreases when $s+c>1$. Unknown sensitivity in the relevant condition limits that calculation. The principle is therefore to state the observation, the target, and the assumptions connecting them. \emph{Unknown} describes the evidence; it is not another experiential state.

\section{Compare definitions against the same constraints}
\subsection*{Criteria for comparison}
A system-level definition should be assessed by five questions. Does it preserve the referent across the documented changes? Does it introduce unnecessary conditions? Does it distinguish the system from its faculties and evidence? Does it use consistent biological boundaries? Does it cover the relevant cases without changing targets?

These criteria evaluate the definition against existing findings. Each claimed exclusion must identify a documented counterexample or explain why the candidate addresses a different target. Comparing definitions does not require collecting a new dataset when the relevant counterexamples are already available.

\subsection*{Faculty-based candidates}
Overt response and particular memory functions fail as necessary conditions of the continuing biological individual in the documented cases. Wakefulness and environmental connectedness likewise describe states rather than organismal existence. Such variables can remain valuable measures of their own targets after their identification with persistence is rejected.

Necessity is distinct from causal support. A neural mechanism may be necessary for a faculty without being identical to the biological individual. Conversely, faculties can affect the individual's survival. A system--faculty distinction is not a claim that causation runs only in one direction.

\subsection*{The strongest rival: an experience-centered definition}
An experience-centered account reserves consciousness for $E$ and separately records biological persistence. It can accept the same dissociations. Chalmers explicitly identifies experience as his hard target \cite{Chalmers1995}. This rival does not fail simply because an organism survives a change in experience: that is compatible with its definition.

SCTC's response is to specify the system-level question that such a definition does not answer by itself: which biological individual continues across changing experiential and functional states? The experience-centered account can answer that question by adding $L$ and $\Gamma$. Where the resulting descriptions make identical observational commitments, the evidence does not establish exclusive victory for either vocabulary.

The defensible comparison is thus target-relative. SCTC offers a minimal definition for organismal persistence and an explicit separation from $E$. Claiming that it is the only possible definition of consciousness would additionally require showing why experience-centered and other viable accounts cannot serve their declared targets. The cited record does not supply that universal exclusion.

\subsection*{Mechanisms and definitions answer different questions}
Global neuronal workspace theory concerns conscious access through broad availability of information; integrated information theory 4.0 connects phenomenal existence to intrinsic causal structure \cite{Dehaene2011,Albantakis2023}. These are additional commitments about access or experience, not merely alternative names for biological persistence.

The Cogitate adversarial study compared preregistered predictions from those theories and reported support and challenges for each \cite{Cogitate2025}. SCTC can organize the biological bearer while asking a mechanistic theory to specify which state it explains. To compete as a mechanism for experience, SCTC would need an additional account of $E$. Its present empirical grounding comes from the existing constraints on system, faculty, and observation.

\section{Boundary consistency and conditions for revision}
\subsection*{An individual is more than viable tissue}
A viable culture, an isolated organ, and an integrated organism are different units. Physiological regulation across parts informs organismal organization \cite{BechtelBich2024}. Application of SCTC requires a stated unit, scale, interval, causal history, and account of integration. Isolated cellular activity is insufficient by itself to classify persistence of the organism.

External support and component replacement must be assessed through those same criteria. A definition cannot classify a trajectory as continuous only because recovery later occurred, then treat recovery as its independently predicted consequence. Use the evidence available at the specified time and state where the boundary remains unresolved.

\subsection*{The scope of the biological assignment}
Applied to all qualifying biological individuals, Eq.~\ref{eq:definition} includes plants and microbes in the system-level category. It does not infer experience, pain, or a point of view in them. Adding a cognitive threshold would change the definition and needs a separate argument.

Nonbiological systems lie outside the declared domain. That restriction does not establish that artificial experience is impossible. Similarly, one organismal trajectory does not establish one stream of experience or preservation of psychological identity. These questions require their own targets and evidence.

\subsection*{What would require a revision?}
The account must be revised if its stated biological criteria cannot classify its intended domain consistently, if a claimed faculty exclusion rests on an incorrectly established dissociation, or if the formulation conflates organismal persistence with an experiential conclusion. A rival system-level account deserves preference if it handles the same evidence and boundaries with fewer unsupported commitments.

A new experiment is needed only when a further claim is not settled by available evidence. If SCTC adds a causal claim about experience, that claim must specify distinct observable consequences and survive comparison with alternatives. Existing dissociations already constrain the current definition; they cannot establish an unspecified future mechanism.

The system-level category alone supplies no inference about clinical awareness, suffering, prognosis, personhood, or moral or legal status. Those are additional empirical or normative questions. Biological boundary analysis does not replace established clinical death criteria.

\section{The hard problem and the system-level result}
The hard problem asks why physical processes are accompanied by experience \cite{Chalmers1995}. In this paper's notation, it concerns $E$. SCTC does not resolve that problem by defining $C_{\mathrm{sys}}$. It distinguishes the problem of experience from the persistence of the individual that may instantiate it.

The distinction also places an obligation on SCTC. Defining its target does not explain the mechanisms that maintain biological organization. Nor does separating $L$ from $E$ remove the need to explain their relation in a comprehensive theory. A definition identifies what an explanation must account for; it is not the explanation itself.

The empirical result used here is that several commonly associated variables do not behave as a single indivisible measure. The analytic result is that failure of an observation channel does not automatically locate failure of the inferred state. The theoretical result is a minimal system-level definition that keeps the continuing biological individual fixed across changes in those faculties.

SCTC is therefore an empirically grounded account of a declared system-level target. It is already constrained by published dissociations and explicit logical requirements. Comparative claims about its superiority must be made against rivals addressing that same target. Experience remains a legitimate scientific target in its own right.

\begin{keybox}{Core conclusion}
\textbf{The individual continues. Its states change. Our evidence can fail. A definition must keep those claims distinct.}
\end{keybox}

The practical reporting rule is equally direct: identify the biological unit, name the state or capacity, state the observed signal, and specify what the evidence warrants. For example, ``No bedside command response was observed; neural command-following evidence was detected; the contents of experience remain unresolved.'' This preserves the finding without turning one measurement into a verdict about every variable.

{\small\textbf{Method and materials.} This paper synthesizes selected published empirical studies and develops a definition and analytic arguments. No new experimental dataset or quantitative reanalysis is reported. The \href{\repoURL}{research repository} contains the current manuscript, LaTeX source, and supporting materials. Numbered citations refer to bibliography entries, not page numbers.\par}

{\small\textbf{AI assistance.} Generative AI tools and coding assistants were used during manuscript preparation for formatting, language editing, LaTeX preparation, layout iteration, literature organization, and adversarial stress-testing of the manuscript. The research questions, theoretical claims, interpretation of the literature, and final manuscript decisions are the author's.\par}

\begin{thebibliography}{99}
\scriptsize
\setlength{\itemsep}{0em}
\setlength{\parsep}{0pt}
\bibitem{Chalmers1995} Chalmers, D. J. (1995). Facing up to the problem of consciousness. Journal of Consciousness Studies, 2(3), 200-219.\newline \href{https://consc.net/papers/facing.html}{https://consc.net/papers/facing.html}

\bibitem{Block1995} Block, N. (1995). On a confusion about a function of consciousness. Behavioral and Brain Sciences, 18, 227-247.\newline \href{https://doi.org/10.1017/S0140525X00038188}{https://doi.org/10.1017/S0140525X00038188}

\bibitem{Bayne2016} Bayne, T., Hohwy, J., \& Owen, A. M. (2016). Are there levels of consciousness? Trends in Cognitive Sciences, 20, 405-413.\newline \href{https://doi.org/10.1016/j.tics.2016.03.009}{https://doi.org/10.1016/j.tics.2016.03.009}

\bibitem{Bodien2024} Bodien, Y. G., et al. (2024). Cognitive motor dissociation in disorders of consciousness. New England Journal of Medicine, 391, 598-608.\newline \href{https://doi.org/10.1056/NEJMoa2400645}{https://doi.org/10.1056/NEJMoa2400645}

\bibitem{Sanders2012} Sanders, R. D., Tononi, G., Laureys, S., \& Sleigh, J. W. (2012). Unresponsiveness \ensuremath{\neq} unconsciousness. Anesthesiology, 116, 946-959.\newline \href{https://doi.org/10.1097/ALN.0b013e318249d0a7}{https://doi.org/10.1097/ALN.0b013e318249d0a7}

\bibitem{Siclari2017} Siclari, F., et al. (2017). The neural correlates of dreaming. Nature Neuroscience, 20, 872-878.\newline \href{https://doi.org/10.1038/nn.4545}{https://doi.org/10.1038/nn.4545}

\bibitem{Milner2005} Milner, B. (2005). The medial temporal-lobe amnesic syndrome. Psychiatric Clinics of North America, 28, 599-611.\newline \href{https://doi.org/10.1016/j.psc.2005.06.002}{https://doi.org/10.1016/j.psc.2005.06.002}

\bibitem{Casey2024} Casey, C. P., et al. (2024). Evaluation of putative signatures of consciousness using specific definitions of responsiveness, connectedness, and consciousness. British Journal of Anaesthesia, 132, 300-311.\newline \href{https://doi.org/10.1016/j.bja.2023.09.031}{https://doi.org/10.1016/j.bja.2023.09.031}

\bibitem{Casali2013} Casali, A. G., et al. (2013). A theoretically based index of consciousness independent of sensory processing and behavior. Science Translational Medicine, 5, 198ra105.\newline \href{https://doi.org/10.1126/scitranslmed.3006294}{https://doi.org/10.1126/scitranslmed.3006294}

\bibitem{Dehaene2011} Dehaene, S., \& Changeux, J.-P. (2011). Experimental and theoretical approaches to conscious processing. Neuron, 70, 200-227.\newline \href{https://doi.org/10.1016/j.neuron.2011.03.018}{https://doi.org/10.1016/j.neuron.2011.03.018}

\bibitem{Albantakis2023} Albantakis, L., et al. (2023). Integrated information theory (IIT) 4.0: formulating the properties of phenomenal existence in physical terms. PLOS Computational Biology, 19, e1011465.\newline \href{https://doi.org/10.1371/journal.pcbi.1011465}{https://doi.org/10.1371/journal.pcbi.1011465}

\bibitem{Cogitate2025} Cogitate Consortium, et al. (2025). Adversarial testing of global neuronal workspace and integrated information theories of consciousness. Nature, 642, 133-142.\newline \href{https://doi.org/10.1038/s41586-025-08888-1}{https://doi.org/10.1038/s41586-025-08888-1}

\bibitem{BechtelBich2024} Bechtel, W., \& Bich, L. (2024). Situating homeostasis in organisms: maintaining organization through time. The Journal of Physiology, 602, 6003-6020.\newline \href{https://doi.org/10.1113/JP286883}{https://doi.org/10.1113/JP286883}

\end{thebibliography}
\end{document}